\begin{ouverture}[radioactivite]{4}{Atomes, noyaux et ions}{Constitution et transformations}{Observer · Modéliser · Expérimenter · Expliquer}
\objectif{Relier l’écriture d’un noyau à sa composition.}
\objectif{Comparer tailles et masses aux échelles atomiques.}
\objectif{Distinguer atome, molécule, cation et anion.}
\prerequis{Proportionnalité, conversions et lecture de graphiques.}
\end{ouverture}
\essentiel
\essSection{1}{Une structure presque vide}
Un atome comporte un noyau positif entouré d'électrons négatifs. Sa taille est de l'ordre de $10^{-10}\unit{m}$, celle d'un noyau de l'ordre de $10^{-15}\unit{m}$. Le rapport des tailles vaut environ $10^5$ : le noyau occupe une très petite partie du volume atomique.
\schema[corrige]{atome}
\begin{definition}[Constituants]Un proton porte $+e$, un électron $-e$, un neutron une charge nulle. $e=\qty{1,602e-19}{C}$. Les masses du proton et du neutron sont proches de \qty{1,67e-27}{kg}, celle de l'électron vaut environ \qty{9,11e-31}{kg}.\end{definition}
\essSection{2}{Identifier un noyau}
L'écriture $\ce{^{A}_{Z}X}$ donne $Z$ protons et $A$ nucléons, donc $A-Z$ neutrons. Le numéro atomique $Z$ définit l'élément chimique. Un atome neutre possède $Z$ électrons. Deux noyaux du même élément peuvent avoir des nombres de neutrons différents.
\begin{exemple}[Un atome de sodium]$\ce{^{23}_{11}Na}$ possède 11 protons, 12 neutrons et, à l'état neutre, 11 électrons. Sa masse est approximativement $23\times1,67\times10^{-27}=\qty{3,84e-26}{kg}$ : la masse des électrons est négligeable à cette précision.\end{exemple}
\essSection{3}{Ions et neutralité de la matière}
Un cation résulte d'une perte d'électrons ; un anion d'un gain d'électrons. Le noyau reste inchangé lors de cette ionisation. $\ce{Na+}$ a 10 électrons, $\ce{Cl-}$ en a 18. Une molécule est un assemblage électriquement neutre d'atomes liés.
\begin{propriete}[Composé ionique]La formule indique les proportions d'ions nécessaires à la neutralité. $\ce{CaCl2}$ associe un ion $\ce{Ca^{2+}}$ à deux ions $\ce{Cl-}$ ; ce solide n'est pas un ensemble de molécules isolées de chlorure de calcium.\end{propriete}
\essSection{4}{Expérimenter et modéliser}
\begin{experience}[Construire un modèle]Avec des jetons, représenter les protons, neutrons et électrons d'un atome, puis de son ion. Dresser deux bilans : nombre de nucléons inchangé, nombre d'électrons modifié. Le modèle représente les nombres et les charges, pas les dimensions réelles ni des orbites planétaires.\end{experience}
\exercices
\rubrique{Prendre confiance}
\begin{blocLivre}{ex-4-1}\begin{exercice}[1]{Lire une écriture nucléaire}
On considère $\ce{^{27}_{13}Al}$.
\begin{questions}
\item Donner les nombres de protons et de neutrons.
\item Donner le nombre d’électrons de l’atome neutre.
\end{questions}
\end{exercice}\end{blocLivre}
\begin{blocLivre}{ex-4-2}\begin{exercice}[1]{Écrire un noyau}
Un noyau d’oxygène contient 8 protons et 10 neutrons.
\begin{questions}
\item Déterminer $Z$ et $A$.
\item Écrire son symbole nucléaire.
\end{questions}
\end{exercice}\end{blocLivre}
\begin{blocLivre}{ex-4-3}\begin{exercice}[1]{Qualifier une entité}
On donne $\ce{Ne}$, $\ce{CO2}$, $\ce{Mg^{2+}}$ et $\ce{F-}$.
\begin{questions}
\item Associer les mots atome, molécule, cation et anion.
\item Quelle entité porte deux charges élémentaires positives ?
\end{questions}
\end{exercice}\end{blocLivre}
\begin{blocLivre}{ex-4-4}\begin{exercice}[1]{Un anion chlorure}
Le chlore possède $Z=17$.
\begin{questions}
\item Combien d’électrons possède \ce{Cl-} ?
\item La formation de cet ion change-t-elle le noyau ?
\end{questions}
\end{exercice}\end{blocLivre}
\rubrique{S’entraîner}
\begin{blocLivre}{ex-4-5}\begin{exercice}[2]{Changer d’échelle}
Un atome mesure environ $10^{-10}\unit{m}$ et son noyau $10^{-15}\unit{m}$.
\schema{atome}
\begin{questions}
\item Calculer le rapport des tailles.
\item Si l’atome avait un diamètre de 100 m, quel serait celui du noyau ?
\end{questions}
\end{exercice}\end{blocLivre}
\begin{blocLivre}{ex-4-6}\begin{exercice}[2]{Une masse atomique}
On assimile les masses des nucléons à \qty{1,67e-27}{kg}. Un atome de carbone possède 12 nucléons.
\begin{questions}
\item Estimer sa masse.
\item Pourquoi la masse des six électrons est-elle négligeable ici ?
\end{questions}
\end{exercice}\end{blocLivre}
\begin{blocLivre}{ex-4-7}\begin{exercice}[2]{Charge du noyau}
Un noyau de magnésium possède 12 protons ; $e=\qty{1,60e-19}{C}$.
\begin{questions}
\item Calculer sa charge.
\item Quelle est la charge totale des électrons de l’atome neutre ?
\end{questions}
\end{exercice}\end{blocLivre}
\begin{blocLivre}{ex-4-8}\begin{exercice}[2]{Neutraliser les charges}
On associe des ions $\ce{Al^{3+}}$ et $\ce{O^{2-}}$.
\begin{questions}
\item Trouver les plus petits nombres d’ions assurant la neutralité.
\item En déduire la formule du composé.
\end{questions}
\end{exercice}\end{blocLivre}
\rubrique{Raisonner et approfondir}
\begin{blocLivre}{ex-4-9}\begin{exercice}[3]{Un ion à identifier}
Une entité possède 11 protons, 12 neutrons et 10 électrons.
\begin{questions}
\item Déterminer sa charge en multiples de $e$.
\item Donner son symbole et le nombre de masse du noyau.
\end{questions}
\end{exercice}\end{blocLivre}
\begin{blocLivre}{ex-4-10}\begin{exercice}[3]{Une erreur de vocabulaire}
Un texte affirme : « Le cation calcium se forme quand le noyau gagne deux protons ».
\begin{questions}
\item Corriger l’explication.
\item Que provoquerait réellement un changement de nombre de protons ?
\end{questions}
\end{exercice}\end{blocLivre}
\begin{blocLivre}{ex-4-11}\begin{exercice}[3]{Masse et charge}
Un ion possède un noyau $\ce{^{24}_{12}Mg}$ et 10 électrons.
\begin{questions}
\item Calculer sa charge avec $e=\qty{1,60e-19}{C}$.
\item Comparer sa masse à celle de l’atome neutre.
\end{questions}
\end{exercice}\end{blocLivre}
\begin{blocLivre}{ex-4-12}\begin{exercice}[3]{Un solide électriquement neutre}
Un échantillon contient $3,0\times10^{20}$ ions $\ce{Ca^{2+}}$ et uniquement des ions chlorure pour compenser.
\begin{questions}
\item Combien d’ions chlorure contient-il ?
\item Calculer la charge totale du solide.
\end{questions}
\end{exercice}\end{blocLivre}

